\subsection{Cascode Mirror}

\subsubsection{Overview}
Constructed from four transistors, It improves
output resistance and accuracy by reducing the
effects of channel length modulation.

\subsubsection{Schematic}
\begin{center}
  \begin{tikzpicture}
  	% Paths, nodes and wires:
  	\node[shape=rectangle, draw={rgb,255:red,255;green,0;blue,0}, line width=1pt, dash pattern={on 4pt off 2pt on 1pt off 2pt on 1pt off 2pt}, minimum width=1.715cm, minimum height=2.195cm] at (6.125, 7.135){};
  	\draw (3.23, 8.04) to[american voltage source, l_={$V_{cc}$}] (1.98, 8.04);
  	\node[nmos, arrowmos, xscale=-1](N1) at (3.5, 2.75){} node[anchor=east] at (N1.text){$M_0$};
  	\node[nmos, arrowmos](N2) at (6, 2.75){} node[anchor=west] at (N2.text){$M_1$};
  	\draw (3.5, 7.54) to[european resistor, l={$G_4$}] (3.5, 6.04);
  	\draw (6, 5.79) |- (6, 5.25);
  	\draw (5.02, 2.75) -- (4.48, 2.75);
  	\draw (3.5, 1.98) -| (3.5, 1.75) -| (6, 1.98);
  	\draw (4, 1) to[american voltage source, l={$V_{ee}$}] (2.75, 1);
  	\draw (4, 1) -| (4.75, 1.75);
  	\node[ground, rotate=-90] at (1.98, 8.04){};
  	\node[ground, rotate=-90] at (2.75, 1){};
  	\draw[dash pattern={on 0.4pt off 0.8pt}, -stealth] (6, 5.5) -- (6.5, 5.5);
  	\node[shape=rectangle, minimum width=0.965cm, minimum height=0.465cm](N3) at (7, 5.46){} node[anchor=center] at (N3.text){$V_{out}$};
  	\draw (6, 7.5) to[european resistor, l={\textcolor{rgb,255:red,255;green,0;blue,0}{$G_5$}}] (6, 6);
  	\draw (6, 5.79) -- (6, 6.02);
  	\draw (3.5, 7.54) |- (3.23, 8.04);
  	\draw (3.5, 8.04) -| (6, 7.5);
  	\draw (10.25, 10.75) -- (10.25, -0);
  	\node[shape=rectangle, minimum width=5.59cm, minimum height=1.09cm] at (5.438, 9.937){} node[anchor=north, align=center, text width=5.202cm, inner sep=6pt] at (5.438, 10.5){Large Signal Schematic\\(Sink)};
  	\node[shape=rectangle, minimum width=3.09cm, minimum height=2.465cm] at (8.687, 7){} node[anchor=north west, align=left, text width=2.702cm, inner sep=6pt] at (7.125, 8.25){\textcolor{rgb,255:red,255;green,0;blue,0}{\footnotesize Exists Only For Analysis. In  Implementation take output directly from V\_\{out\}}};
  	\node[shape=rectangle, minimum width=0.965cm, minimum height=0.465cm](N4) at (7, 3.71){} node[anchor=center] at (N4.text){$V_{y}$};
  	\draw[dash pattern={on 0.4pt off 0.8pt}, -stealth] (6, 3.71) |- (6.5, 3.71);
  	\node[shape=rectangle, minimum width=0.965cm, minimum height=0.465cm](N5) at (4.75, 4.25){} node[anchor=center] at (N5.text){$I_{G_1}$};
  	\draw (17.625, 2.75) to[american current source, l_={$g_{m_1}v_{gs}$}] (17.625, 1.25);
  	\draw (19.125, 2.75) to[european resistor, l_={$g_{o_1}$}] (19.125, 1.25);
  	\draw (17.625, 1.25) -- (19.125, 1.25);
  	\draw (19.125, 2.75) -- (17.625, 2.75);
  	\draw (17.625, 1.25) -- (16.375, 1.25);
  	\draw (16.375, 2.75) -- (15.625, 2.75);
  	\node[shape=rectangle, minimum width=0.465cm, minimum height=0.465cm] at (16.375, 2.5){} node[anchor=center, align=center, text width=0.077cm, inner sep=6pt] at (16.375, 2.5){\footnotesize G};
  	\node[shape=rectangle, minimum width=0.465cm, minimum height=0.465cm] at (16.375, 1.5){} node[anchor=center, align=center, text width=0.077cm, inner sep=6pt] at (16.375, 1.5){\footnotesize S};
  	\node[shape=rectangle, minimum width=0.465cm, minimum height=0.465cm] at (19.625, 2.5){} node[anchor=center, align=center, text width=0.077cm, inner sep=6pt] at (19.625, 2.5){\footnotesize D};
  	\draw (15.625, 2.75) -- (14.875, 2.75);
  	\draw (13.375, 2.75) to[american current source, l={$g_{m_0}v_{gs}$}] (13.375, 1.25);
  	\draw (11.875, 2.75) to[european resistor, l={$g_{o_0}$}] (11.875, 1.25);
  	\draw (13.375, 2.75) -- (11.875, 2.75);
  	\draw (13.375, 1.25) -- (11.875, 1.25);
  	\draw (11.875, 2.75) -| (11.875, 4.25);
  	\draw (14.875, 1.25) -- (13.375, 1.25);
  	\node[shape=rectangle, minimum width=0.465cm, minimum height=0.465cm] at (14.875, 2.5){} node[anchor=center, align=center, text width=0.077cm, inner sep=6pt] at (14.875, 2.5){\footnotesize G};
  	\node[shape=rectangle, minimum width=0.465cm, minimum height=0.465cm] at (14.875, 1.5){} node[anchor=center, align=center, text width=0.077cm, inner sep=6pt] at (14.875, 1.5){\footnotesize S};
  	\node[shape=rectangle, minimum width=0.465cm, minimum height=0.465cm] at (11.375, 2.5){} node[anchor=center, align=center, text width=0.077cm, inner sep=6pt] at (11.375, 2.5){\footnotesize D};
  	\draw (11.875, 8.562) to[european resistor, l={$G_4$}] (11.875, 7.062);
  	\draw (11.875, 7.062) -| (11.875, 6.562);
  	\draw (11.875, 8.562) -| (11.875, 9.062) -- (12.625, 9.062) -| (12.625, 8.812);
  	\draw (19.125, 7.75) to[european resistor, l={\textcolor{rgb,255:red,255;green,0;blue,0}{$G_5$}}] (19.125, 6.25);
  	\draw (19.125, 6.25) -| (19.125, 5.75);
  	\draw (19.125, 7.75) -| (19.125, 8.25) |- (18.375, 8.25) -| (18.375, 8);
  	\node[shape=rectangle, minimum width=0.715cm, minimum height=0.465cm](N6) at (20, 3.5){} node[anchor=center] at (N6.text){$V_{y}$};
  	\node[shape=rectangle, minimum width=0.965cm, minimum height=0.465cm](N7) at (20.125, 5.75){} node[anchor=center] at (N7.text){$V_{out}$};
  	\node[ground] at (14.875, 1.25){};
  	\node[ground] at (16.375, 1.25){};
  	\node[ground] at (12.625, 8.812){};
  	\node[ground] at (18.375, 8){};
  	\node[shape=rectangle, minimum width=0.965cm, minimum height=0.465cm](N8) at (16.625, 3.563){} node[anchor=center] at (N8.text){$I_{G_1}$};
  	\node[shape=rectangle, minimum width=4.84cm, minimum height=1.09cm] at (15.25, 9.937){} node[anchor=north, align=center, text width=4.452cm, inner sep=6pt] at (15.25, 10.5){Small Signal Schematic\\(Sink)};
  	\draw (6, 3.52) -| (6, 3.71);
  	\node[nmos, arrowmos](N9) at (6, 4.48){} node[anchor=west] at (N9.text){$M_2$};
  	\node[shape=rectangle, minimum width=0.965cm, minimum height=0.465cm](N10) at (2.5, 3.71){} node[anchor=center] at (N10.text){$V_{x}$};
  	\draw[dash pattern={on 0.4pt off 0.8pt}, -stealth] (3.5, 3.71) |- (3, 3.71);
  	\draw (17.625, 5.75) to[american current source, l_={$g_{m_2}v_{gs}$}] (17.625, 4.25);
  	\draw (19.125, 5.75) to[european resistor, l_={$g_{o_2}$}] (19.125, 4.25);
  	\draw (17.625, 4.25) -- (19.125, 4.25);
  	\draw (19.125, 5.75) -- (17.625, 5.75);
  	\draw (17.625, 4.25) -- (16.375, 4.25);
  	\draw (16.375, 5.75) -- (15.625, 5.75);
  	\node[shape=rectangle, minimum width=0.465cm, minimum height=0.465cm] at (16.375, 5.5){} node[anchor=center, align=center, text width=0.077cm, inner sep=6pt] at (16.375, 5.5){\footnotesize G};
  	\node[shape=rectangle, minimum width=0.465cm, minimum height=0.465cm] at (16.375, 4.5){} node[anchor=center, align=center, text width=0.077cm, inner sep=6pt] at (16.375, 4.5){\footnotesize S};
  	\draw (19.125, 4.25) -- (19.125, 2.75);
  	\draw[dash pattern={on 0.4pt off 0.8pt}, -stealth] (19.125, 3.5) -- (19.625, 3.5);
  	\node[shape=rectangle, minimum width=0.715cm, minimum height=0.465cm](N11) at (11, 3.563){} node[anchor=center] at (N11.text){$V_{x}$};
  	\draw[dash pattern={on 0.4pt off 0.8pt}, -stealth] (11.875, 3.563) -- (11.375, 3.563);
  	\node[nmos, arrowmos, xscale=-1](N12) at (3.5, 4.48){} node[anchor=east] at (N12.text){$M_3$};
  	\draw (3.5, 3.71) -| (3.5, 3.52);
  	\draw (4.48, 4.48) -- (5.02, 4.48);
  	\draw (4.75, 4.5) |- (3.5, 5.5);
  	\draw (3.5, 5.25) -| (3.5, 5.5);
  	\draw (3.5, 6.04) -| (3.5, 5.5);
  	\node[shape=rectangle, minimum width=0.965cm, minimum height=0.465cm](N13) at (2.5, 5.5){} node[anchor=center] at (N13.text){$V_{z}$};
  	\draw[dash pattern={on 0.4pt off 0.8pt}, -stealth] (3.5, 5.5) |- (3, 5.5);
  	\node[shape=rectangle, minimum width=0.965cm, minimum height=0.465cm](N14) at (4.75, 6){} node[anchor=center] at (N14.text){$I_{G_2}$};
  	\draw (15.625, 5.75) -- (14.875, 5.75);
  	\draw (13.375, 5.75) to[american current source, l={$g_{m_3}v_{gs}$}] (13.375, 4.25);
  	\draw (11.875, 5.75) to[european resistor, l={$g_{o_3}$}] (11.875, 4.25);
  	\draw (13.375, 5.75) -- (11.875, 5.75);
  	\draw (13.375, 4.25) -- (11.875, 4.25);
  	\draw (14.875, 4.25) -- (13.375, 4.25);
  	\node[shape=rectangle, minimum width=0.465cm, minimum height=0.465cm] at (14.875, 5.5){} node[anchor=center, align=center, text width=0.077cm, inner sep=6pt] at (14.875, 5.5){\footnotesize G};
  	\node[shape=rectangle, minimum width=0.465cm, minimum height=0.465cm] at (14.875, 4.5){} node[anchor=center, align=center, text width=0.077cm, inner sep=6pt] at (14.875, 4.5){\footnotesize S};
  	\node[shape=rectangle, minimum width=0.465cm, minimum height=0.465cm] at (11.375, 5.5){} node[anchor=center, align=center, text width=0.077cm, inner sep=6pt] at (11.375, 5.5){\footnotesize D};
  	\draw (15.625, 5.75) |- (11.875, 6.562) |- (11.875, 5.75);
  	\node[shape=rectangle, minimum width=0.715cm, minimum height=0.465cm](N15) at (11, 6.562){} node[anchor=center] at (N15.text){$V_{z}$};
  	\draw[dash pattern={on 0.4pt off 0.8pt}, -stealth] (11.875, 6.562) -- (11.375, 6.562);
  	\node[shape=rectangle, minimum width=0.965cm, minimum height=0.465cm](N16) at (16.625, 6.562){} node[anchor=center] at (N16.text){$I_{G_2}$};
  	\draw[dash pattern={on 0.4pt off 0.8pt}, -stealth] (19.125, 5.75) -- (19.625, 5.75);
  	\draw[dash pattern={on 0.4pt off 0.8pt}, -stealth] (4.75, 5.5) -- (4.75, 5.75);
  	\draw[dash pattern={on 0.4pt off 0.8pt}, -stealth] (4.75, 3.71) -- (4.75, 4);
  	\draw[dash pattern={on 0.4pt off 0.8pt}, -stealth] (15.625, 6.562) -- (16.125, 6.562);
  	\draw[dash pattern={on 0.4pt off 0.8pt}, -stealth] (15.625, 3.563) -- (16.125, 3.563);
  	\node[shape=rectangle, draw={rgb,255:red,255;green,0;blue,0}, line width=1pt, dash pattern={on 4pt off 2pt on 1pt off 2pt on 1pt off 2pt}, minimum width=1.215cm, minimum height=2.195cm] at (19.375, 7.385){};
  	\draw (3.5, 3.71) -| (4.75, 2.75);
  	\draw (15.625, 2.75) |- (11.875, 3.563);
  \end{tikzpicture}
\end{center}

\subsubsection{Transfer Function}
Here's a table of the voltages at each terminal of the
transistors:\\
\begin{tabularx}{\linewidth}{|X|X|X|X|X|X|} \hline
  Transistor & Gate (G) & Source (S) & Drain (D) & $V_{gs}$ & $V_{ds}$ \\ \hline
  $M_0$ & $V_{x}$ & $V_{ee}$ & $V_{x}$ & $V_{x} - V_{ee}$ & $V_{x} - V_{ee}$ \\ \hline
  $M_1$ & $V_{x}$ & $V_{ee}$ & $V_{y}$ & $V_{x} - V_{ee}$ & $V_{y} - V_{ee}$ \\ \hline
  $M_2$ & $V_{z}$ & $V_{y}$ & $V_{out}$ & $V_{z} - V_{y}$ & $V_{out} - V_{y}$ \\ \hline
  $M_3$ & $V_{z}$ & $V_{x}$ & $V_{z}$ & $V_{z} - V_{x}$ & $V_{z} - V_{x}$ \\ \hline
\end{tabularx}\\

\begin{enumerate}
  \item \textbf{Large Signal Analysis}:\\
    Applying KCL at $V_{out}$:
    \begin{align*}
      (V_{out} - V_{cc})G_5 + I_{M_2} &= 0 \\
      V_{out}G_5 - V_{cc}G_5 + I_{M_2} &= 0 \\
      V_{out} = V_{cc} - \frac{I_{M_2}}{G_5} \tag{1}
    \end{align*}
    Applying KCL at $V_{x}$:
    \begin{align*}
      I_{M_0} + I_{G_1} - I_{M_3} &= 0 \\
      I_{M_0} + I_{G_1} = I_{M_3} \\
      I_{M_0} &\approx I_{M_3} \quad \text{[assuming } I_{G_1} \approx 0 \text{ by MOS physics]} \tag{2}
    \end{align*}
    Applying KCL at $V_{y}$:
    \begin{align*}
      I_{M_1} - I_{M_2} &= 0 \\
      I_{M_1} = I_{M_2} \tag{3}
    \end{align*}
    Applying KCL at $V_{z}$:
    \begin{align*}
      (V_{z} - V_{cc})G_4 + I_{M_3} + I_{G_2} &= 0 \\
      V_{z}G_4 - V_{cc}G_4 + I_{M_3} + I_{G_2} &= 0 \\
      V_{z}G_4 &= V_{cc}G_4 - (I_{M_3} + I_{G_2}) \\
      V_{z} &= V_{cc} - \frac{(I_{M_3} + I_{G_2})}{G_4} \tag{4}
    \end{align*}
    Using the transistor equations for $I_{M_0}$, $I_{M_1}$, $I_{M_2}$, and $I_{M_3}$:
    \begin{align*}
      I_{M_0} &= k_{n_0} (V_{gs_0} - V_{th})^2 (1 + \lambda V_{ds_0}) \\
              &= k_{n_0} (V_{x} - V_{ee} - V_{th})^2 (1 + \lambda (V_{x} - V_{ee})) \\
              &= k_{n_0} (V_{x} - V_{ee} - V_{th})^2 \quad \text{[assuming } \lambda \approx 0 \text{]} \tag{5} \\
      I_{M_1} &= k_{n_1} (V_{gs_1} - V_{th})^2 (1 + \lambda V_{ds_1}) \\
              &= k_{n_1} (V_{x} - V_{ee} - V_{th})^2 (1 + \lambda (V_{y} - V_{ee})) \\
              &= k_{n_1} (V_{x} - V_{ee} - V_{th})^2 \quad \text{[assuming } \lambda \approx 0 \text{]} \tag{6} \\
      I_{M_2} &= k_{n_2} (V_{gs_2} - V_{th})^2 (1 + \lambda V_{ds_2}) \\
              &= k_{n_2} (V_{z} - V_{y} - V_{th})^2 (1 + \lambda (V_{out} - V_{y})) \\
              &= k_{n_2} (V_{z} - V_{y} - V_{th})^2 \quad \text{[assuming } \lambda \approx 0 \text{]} \tag{7} \\
      I_{M_3} &= k_{n_3} (V_{gs_3} - V_{th})^2 (1 + \lambda V_{ds_3}) \\
              &= k_{n_3} (V_{z} - V_{x} - V_{th})^2 (1 + \lambda (V_{z} - V_{x})) \\
              &= k_{n_3} (V_{z} - V_{x} - V_{th})^2 \quad \text{[assuming } \lambda \approx 0 \text{]} \tag{8}
    \end{align*}
    It can be seen that,
    \begin{align*}
      I_{in} = I_{M_0} &= I_{M_3} \quad \text{[from (2)]} \\
      I_{out} = I_{M_1} &= I_{M_2} \quad \text{[from (3)]} \\
      I_{M_0} &= I_{M_1} \quad \text{[from (5) and (6) and assuming } k_{n_0} = k_{n_1} \text{]} \\
      \therefore, I_{in} &= I_{out} \\
    \end{align*}
    Thus,
    \begin{align*}
      \boxed{\frac{I_{out}}{I_{in}} \approx \frac{I_{M_2}}{I_{M_3}} \approx \frac{I_{M_1}}{I_{M_0}}
                                                                    \approx \frac{k_{n_1}}{k_{n_0}} = \frac{\frac{\mu_n C_{ox} W_1}{L_1}}{\frac{\mu_n C_{ox} W_0}{L_0}} = \frac{\frac{W_1}{L_1}}{\frac{W_0}{L_0}} = \frac{W_1 L_0}{L_1 W_0}
                                                                  = \frac{W_2 L_3}{L_2 W_3}} \\
    \end{align*}
  \item \textbf{Small Signal Analysis}:\\
    Applying KCL at $V_{out}$:
    \begin{align*}
      V_{out}G_5 + (V_{out} - V_{y})g_{o_2} + (V_{z} - V_{y})g_{m_2} &= 0 \\
      V_{out}G_5 + V_{out}g_{o_2} - V_{y}g_{o_2} + V_{z}g_{m_2} - V_{y}g_{m_2} &= 0 \\
      V_{out}(G_5 + g_{o_2}) - V_{y}(g_{o_2} + g_{m_2}) + V_{z}g_{m_2} &= 0 \tag{9}
    \end{align*}
    Applying KCL at $V_{x}$:
    \begin{align*}
      V_{x}g_{o_0} + I_{G_1} + V_{x}g_{m_0} + (V_{x} - V_{z})g_{o_3} - (V_{z} - V_{x})g_{m_3} &= 0 \\
      V_{x}g_{o_0} + V_{x}g_{m_0} + (V_{x} - V_{z})g_{o_3} - (V_{z} - V_{x})g_{m_3} &= 0 \quad \text{[as } I_{G} \approx 0 \text{ by MOS physics]} \\
      V_{x}g_{o_0} + V_{x}g_{m_0} + V_{x}g_{o_3} - V_{z}g_{o_3} - V_{z}g_{m_3} + V_{x}g_{m_3} &= 0 \\
      V_{x}(g_{o_0} + g_{m_0} + g_{o_3} + g_{m_3}) - V_{z}(g_{o_3} + g_{m_3}) &= 0 \tag{10}
    \end{align*}
    Applying KCL at $V_{y}$:
    \begin{align*}
      V_{y}g_{o_1} + V_{x}g_{m_1} + (V_{y} - V_{out})g_{o_2} - (V_{z} - V_{y})g_{m_2} &= 0 \\
      V_{y}g_{o_1} + V_{x}g_{m_1} + V_{y}g_{o_2} - V_{out}g_{o_2} - V_{z}g_{m_2} + V_{y}g_{m_2} &= 0 \\
      V_{y}(g_{o_1} + g_{o_2} + g_{m_2}) - V_{out}g_{o_2} - V_{z}g_{m_2} + V_{x}g_{m_1} &= 0 \tag{11}
    \end{align*}
    Applying KCL at $V_{z}$:
    \begin{align*}
      V_{z}G_4 + (V_{z} - V_{x})g_{o_3} + (V_{z} - V_{x})g_{m_3} + I_{G_2} &= 0 \\
      V_{z}G_4 + (V_{z} - V_{x})g_{o_3} + (V_{z} - V_{x})g_{m_3} &= 0 \\quad \text{[as } I_{G} \approx 0 \text{ by MOS physics]} \\
      V_{z}G_4 + V_{z}g_{o_3} - V_{x}g_{o_3} + V_{z}g_{m_3} - V_{x}g_{m_3} &= 0 \\
      V_{z}(G_4 + g_{o_3} + g_{m_3}) - V_{x}(g_{o_3} + g_{m_3}) &= 0 \tag{12}
    \end{align*}
    Adding voltage source $V_s$ at $V_z$ that supplies
    a current $I_s$ into $V_z$, we get:
    \begin{align*}
      V_{z}(G_4 + g_{o_3} + g_{m_3}) - V_{x}(g_{o_3} + g_{m_3}) - I_s &= 0 \tag{13}
    \end{align*}
    Fifth equation arises from Voltage source:
    \begin{align*}
      V_{z} = V_s \tag{14}
    \end{align*}
    Equations (9), (10), (11), (13), and (14) can be
    solved analytically. It is cumbersome so, I've
    chosen to give the final result directly:
    \begin{align*}
      \frac{V_{out}}{V_{s}} \quad \text{where,} \\
      V_{out} = - g_{m_0}g_{m_2}g_{o_1} - g_{m_1}g_{m_2}(g_{m_3} + g_{o_3}) - g_{m_1}g_{o_2}(g_{m_3} + g_{o_3}) - g_{m_2}g_{o_1}(g_{m_3} + (g_{o_0} + g_{o_3})) \\
      V_{s} = G_{5}[(g_{m_0}(g_{m_2} + g_{o_1} + g_{o_2}) + g_{m_2}(g_{m_3} + g_{o_0} + g_{o_3}) + g_{m_3}(g_{o_1} + g_{o_2})) + g_{o_0}(g_{o_1} + g_{o_2}) \\
              + (g_{o_3}(g_{o_1} + g_{o_2}))] + g_{o_1}g_{o_2}(g_{m_0} + g_{m_3} + g_{o_0} + g_{o_3})
    \end{align*}
\end{enumerate}


\subsubsection{Analysis}
\begin{enumerate}
  \item \textbf{Output Load (Max)}: Maximum possible
    load that can be connected at the output before
    the transistor $M_2$ goes out of saturation can
    be calculated as follows:
    \begin{align*}
      R_{out, max} &= \frac{V_{L, max}}{I_{M_2}} \quad \text{[where, } V_{ee} \leq V_{L, max} \leq V_{cc} \text{]}
    \end{align*}
  \item \textbf{Impedances}:
    \paragraph{At Input:}Applying a test current $I_{test}$
      at the input node $V_{z}$, from equation (12):
      \begin{align*}
        V_{z}(G_4 + g_{o_3} + g_{m_3}) - V_{x}(g_{o_3} + g_{m_3}) - I_{test} &= 0 \\
        V_{z}(G_4 + g_{o_3} + g_{m_3}) - V_{x}(g_{o_3} + g_{m_3}) &= I_{test} \tag{15}
      \end{align*}
      Solving equations (9), (10), (11), (15) simultaneously for $\frac{V_{z}}{I_{test}}$
      gives the input impedance. It is cumbersome so, I've chosen to give the final
      result directly from Octave:
      \begin{align*}
        Z_{in} = \frac{V_{z}}{I_{test}} = \frac{g_{m_0} + g_{m_3} + g_{o_0} + g_{o_3}}{G_{4}(g_{m_3} + g_{o_3}) + (g_{m_0} + g_{o_0})(G_{4} + g_{m_3} + g_{o_3})}
      \end{align*}
    \paragraph{At Output:}Applying a test current $I_{test}$
      at the output node $V_{out}$, from equation (9):
      \begin{align*}
        V_{out}(G_5 + g_{o_2}) - V_{y}(g_{o_2} + g_{m_2}) + V_{z}g_{m_2} - I_{test} &= 0 \\
        V_{out}(G_5 + g_{o_2}) - V_{y}(g_{o_2} + g_{m_2}) + V_{z}g_{m_2} &= I_{test} \tag{16}
      \end{align*}
      Solving equations (10), (11), (12) and (16) simultaneously for $\frac{V_{out}}{I_{test}}$
      gives the output impedance. It is cumbersome so, I've chosen to give the final
      result directly from Octave:
      \begin{align*}
        Z_{out} = \frac{V_{out}}{I_{test}} = \frac{g_{m_2} + g_{o_1} + g_{o_2}}{G_{5}(g_{m_2} + g_{o_1} + g_{o_2}) + g_{o_1} g_{o_2}}
      \end{align*}
\end{enumerate}


\subsubsection{Design}
Large Signal Analysis can be used to set the bias
currents and voltages in this circuit.
\begin{enumerate}
  \item Choose appropriate transistor dimensions
    ($W$ and $L$) for transistors $M_0$ and $M_1$
    to achieve the desired current gain.
    \begin{align*}
      \frac{I_{out}}{I_{in}} &= \frac{W_1 L_0}{L_1 W_0} = \frac{W_2 L_3}{L_2 W_3} \\
    \end{align*}
  \item Choose a suitable overdrive voltage value
    ($V_{ov}$) for computing the bias voltages at
    nodes $V_{x}$, $V_{y}$, and $V_{z}$.
    \begin{align*}
      V_{ov_0} &= V_{gs_0} - V_{th} \\
      V_{ov_0} &= V_{x} - V_{ee} - V_{th} \\
      V_{x} &= V_{ov_0} + V_{ee} + V_{th} \\
      \\
      V_{ov_3} &= V_{gs_3} - V_{th} \\
      V_{ov_3} &= V_{z} - V_{x} - V_{th} \\
      V_{z} &= V_{ov_3} + V_{x} + V_{th} \\
      \\
      V_{ov_2} &= V_{gs_2} - V_{th} \\
      V_{ov_2} &= V_{z} - V_{y} - V_{th} \\
      V_{y} &= V_{z} - V_{ov_2} - V_{th}
    \end{align*}
  \item Using the chosen $V_{ov}$, compute $W_0$
    to set the required bias current $I_{in}$.
    \begin{align*}
      I_{in} &\approx I_{M_0} = \frac{1}{2} \mu_n C_{ox} \frac{W_0}{L_0} V_{ov_0}^2 \\
             &\approx I_{M_3} = \frac{1}{2} \mu_n C_{ox} \frac{W_3}{L_3} V_{ov_3}^2 \\
      \therefore W_0 &= \frac{2 I_{in} L_0}{\mu_n C_{ox} V_{ov_0}^2} \\
                 W_3 &= \frac{2 I_{in} L_3}{\mu_n C_{ox} V_{ov_3}^2}
    \end{align*}
    $W_1$ can then be computed using the current
    gain relation from step (a). Similarly, $W_2$
    can be computed using the current gain relation
    from step (a).
  \item Compute the required value of $G_4$ (or
    $R_4$) to ensure that the voltage at node
    $V_{z}$ is as calculated in step (b).
    \begin{align*}
      V_{z} &= V_{cc} - \frac{(I_{M_3} + I_{G_2})}{G_4} \\
      V_{z} &\approx V_{cc} - \frac{I_{M_3}}{G_4} \quad \text{[assuming } I_{G_2} \approx 0 \text{ by MOS physics]} \\
      V_{z} - V_{cc} &\approx - \frac{I_{M_3}}{G_4} \\
      \frac{(I_{M_3})}{(V_{cc} - V_{z})} &\approx G_4 \\
      \frac{(I_{in})}{(V_{cc} - V_{z})} &\approx G_4 \quad \text{[} I_{in} = I_{M_3} = I_{M_0} \text{]} \\
      \therefore R_4 &\approx \frac{V_{cc} - V_{z}}{I_{in}}
    \end{align*}
\end{enumerate}

\subsubsection{Question}
Design a cascode current mirror with a current
gain of 5. The input current is $I_{in} = 10 \mu$A.
Assume the following specificiatons:
\begin{enumerate}
  \item $V_{cc} = +5$ V, $V_{ee} = 0$ V
  \item $V_{th} = 0.7$ V
  \item $\mu_n C_{ox} = 100 \mu A/V^2$
  \item $\lambda = 0.02$ V$^{-1}$
  \item $L_{min} = 1 \mu m$
  \item $V_{L, max} = V_{cc}$
  \item $V_{ov} = 0.2$ V for all transistors
\end{enumerate}
Find the following:
\begin{enumerate}
  \item Transistor dimensions ($W$ and $L$) for all
    transistors.
  \item Values of resistor $R_4$ required to set
    the bias current.
  \item Output voltage range and maximum output
    load that can be connected at the output.
  \item Small signal output resistance of the
    current mirror. Assume a load resistance
    $R_L = 100 k\Omega$ is connected at the
    output.
\end{enumerate}
\paragraph{Solution:}
\begin{enumerate}
  \item Transistor dimensions:
    \begin{align*}
      \frac{I_{out}}{I_{in}} &= \frac{W_1 L_0}{L_1 W_0} = \frac{W_2 L_3}{L_2 W_3} = 5 \\
      \\
      V_{x} &= V_{ov_0} + V_{ee} + V_{th} = 0.2 + 0 + 0.7 = 0.9 \text{ V} \\
      V_{z} &= V_{ov_3} + V_{x} + V_{th} = 0.2 + 0.9 + 0.7 = 1.8 \text{ V} \\
      V_{y} &= V_{z} - V_{ov_2} - V_{th} = 1.8 - 0.2 - 0.7 = 0.9 \text{ V} \\
      \\
      W_0 &= \frac{2 I_{in} L_0}{\mu_n C_{ox} V_{ov_0}^2} = \frac{2 \times 10 \times 10^{-6} \times 1 \times 10^{-6}}{100 \times 10^{-6} \times (0.2)^2} = 5 \mu m \\
      W_3 &= \frac{2 I_{in} L_3}{\mu_n C_{ox} V_{ov_3}^2} = \frac{2 \times 10 \times 10^{-6} \times 1 \times 10^{-6}}{100 \times 10^{-6} \times (0.2)^2} = 5 \mu m \\
      W_1 &= 5 \times \frac{W_0}{L_0} \times L_1 = 5 \times \frac{5}{1} \times 1 = 25 \mu m \\
      W_2 &= 5 \times \frac{W_3}{L_3} \times L_2 = 5 \times \frac{5}{1} \times 1 = 25 \mu m
    \end{align*}
  \item Value of resistor $R_4$:
    \begin{align*}
      R_4 &\approx \frac{V_{cc} - V_{z}}
                      {I_{in}} = \frac{5 - 1.8}{10 \times 10^{-6}} = 320 k\Omega
    \end{align*}
  \item Output voltage range and maximum output load:
    \begin{align*}
      V_{out, min} &\approx V_{y} + V_{ov_2} = 0.9 + 0.2 = 1.1 \text{ V} \\
      V_{out, max} &\approx V_{cc} = 5 \text{ V} \\
      R_{out, max} &= \frac{V_{L, max}}{I_{M_2}} = \frac{5 - 0}{10 \times 10^{-6}} = 500 k\Omega
    \end{align*}
  \item Small signal output resistance:
    \begin{align*}
      g_{m_2} &= \frac{2 I_{M_2}}{V_{ov_2}} = \frac{2 \times 10 \times 10^{-6}}{0.2} = 0.0001 \text{ S} \\
      g_{o_1} &= I_{M_1} \lambda = 10 \times 10^{-6} \times 0.02 = 0.2 \times 10^{-6} \text{ S} \\
      g_{o_2} &= I_{M_2} \lambda = 10 \times 10^{-6} \times 0.02 = 0.2 \times 10^{-6} \text{ S} \\
      G_5 &= \frac{1}{R_L} = \frac{1}{100 k\Omega} = 10^{-5} \text{ S} \\
      \\
      Z_{out} &= \frac{g_{m_2} + g_{o_1} + g_{o_2}}{G_{5}(g_{m_2} + g_{o_1} + g_{o_2}) + g_{o_1} g_{o_2}} \\
              &= \frac{0.0001 + 0.2 \times 10^{-6} + 0.2 \times 10^{-6}}{10^{-5}(0.0001 + 0.2 \times 10^{-6} + 0.2 \times 10^{-6}) + (0.2 \times 10^{-6})(0.2 \times 10^{-6})} \\
              &\approx 9.99 M\Omega
    \end{align*}
\end{enumerate}

\subsubsection{Code}
\begin{lstlisting}[language=Octave]
  clear; clc;
  pkg load symbolic;

  syms Vx Vy Vz Vout Vs Is Itest
  syms G5 G4 go0 gm0 go1 gm1 go2 gm2 go3 gm3

  % system kcl
  eq1 = Vout*(G5 + go2) - Vy*(go2 + gm2) + Vz*gm2 == 0; % KCL at Vout
  eq2 = Vx*(go0 + gm0 + go3 + gm3) - Vz*(go3 + gm3) == 0; % KCL at Vx
  eq3 = Vy*(go1 + go2 + gm2) - Vout*go2 - Vz*gm2 + Vx*gm1 == 0; % KCL at Vy
  eq4 = Vz*(G4 + go3 + gm3) - Vx*(go3 + gm3) == 0; % KCL at Vz

  % voltage transfer function: Vs forced at Vz
  eq4_vt = Vz*(G4 + go3 + gm3) - Vx*(go3 + gm3) - Is == 0;
  eq5_vt = Vz == Vs;
  sol_vt = solve([eq1, eq2, eq3, eq4_vt, eq5_vt], [Vout, Vx, Vy, Vz, Vs]);
  volt_transfer = simplify(sol_vt.Vout/sol_vt.Vs);
  % volt_transfer = collect(volt_transfer, G5);
  % volt_transfer = subs(volt_transfer, 'G4', 0);
  % volt_transfer = subs(volt_transfer, 'G5', 0);
  printf('=== Voltage Transfer Results ===\n');
  pretty(volt_transfer)
  printf('\n');

  % current transfer function: Is injected into Vz
  sol_ct = solve([eq1, eq2, eq3, eq4_vt], [Vout, Vx, Vy, Vz]);
  Iout = sol_ct.Vout * G5;
  curr_transfer = simplify(Iout/Is);
  % curr_transfer = collect(curr_transfer, G4*G5);
  printf('=== Current Transfer Results ===\n');
  pretty(curr_transfer)
  printf('\n');

  % input impedence: Itest injected into Vz
  eq4_zin = Vz*(G4 + go3 + gm3) - Vx*(go3 + gm3) - Itest == 0;
  sol_zin = solve([eq1, eq2, eq3, eq4_zin], [Vout, Vx, Vy, Vz]);
  zin = simplify(sol_zin.Vz/Itest);
  % zin = collect(zin, G4*G5);
  % zin = collect(zin, G4);
  % zin = subs(zin, 'G5', 0);
  printf('=== Input Impedence Results ===\n');
  pretty(zin);
  printf('\n');

  % output impedence: Itest injected into Vout
  eq1_zout = Vout*(G5 + go2) - Vy*(go2 + gm2) + Vz*gm2 - Itest == 0;
  sol_zout = solve([eq1_zout, eq2, eq3, eq4], [Vout, Vx, Vy, Vz]);
  zout = simplify(sol_zout.Vout/Itest);
  % zout = collect(zout, G4*G5);
  % zout = subs(zout, 'G4', 0);
  % zout = collect(zout, G5);
  printf('=== Output Impedence Results ===\n');
  pretty(zout)
  printf('\n');
\end{lstlisting}
